

\lussec 7. Calculation of correlation functions

In numerical lattice QCD,
correlation functions involving fields 
at non-zero flow time can be computed
following the steps usually taken in the 
case of ordinary hadronic correlation functions.
The fact that the flow equations must be solved at 
some point nevertheless represents a complication
that needs to be carefully considered.
For illustration, the details are worked out in this section
for two representative cases.


\lussubsec 7.1 Pseudo-scalar two-point function

In practice one is interested in the
correlation function
\equation{
  \langle \Pax^{rs}(x)\Pax_t^{sr}(y)\rangle=
  -\sum_{v,w}\langle\tr\{
  [K(t,y;0,v)S(v,x)_{ss}]^{\dagger}K(t,y;0,w)S(w,x)_{rr}\}\rangle
  \enum
}
at vanishing spatial momentum and computes
its average over a lattice $\Gamma$ of source points,
using random source fields [\cite{MichaelPeisa}],
in order to reduce the statistical errors.

If $x$ is taken to be
the source point, the averaging amounts to the substitution
\equation{
  P^{rs}(x)\to {1\over N_{\Gamma}N_s}
  \sum_{k=1}^{N_s}
  (\psibar_r,\eta_k)(\eta_k,\dirac{5}\psi_s),
  \enum
}
where $N_{\Gamma}$ is the number of points in $\Gamma$ and
\equation{
  \eta_k(x),\quad k=1,\ldots,N_s,
  \enum
}
are randomly chosen complex spinor fields on $\Gamma$ with mean zero
and variance
\equation{
  \langle\eta_k(v)\eta_l(w)^{\dagger}\rangle=
  \cases{\delta_{kl}\delta_{vw} & if $v,w\in\Gamma$,\cr
  \noalign{\vskip1.0ex}
  0 & otherwise.\cr}
  \enum
}
In a QCD simulation,
the computation of the 
averaged correlation function at zero spatial momentum,
\equation{
  {1\over N_{\Gamma}}\sum_{x\in\Gamma}\sum_{\vec{y}}
  \langle \Pax^{rs}(x)\Pax_t^{sr}(y)\rangle=
  -{1\over N_{\Gamma}N_s}\sum_{k=1}^{N_s}
  \sum_{\vec{y}}\langle\phi_{k,s}(t,y)^{\dagger}\phi_{k,r}(t,y)\rangle,
  \enum
}
then amounts to calculating the functions
\equation{
  \phi_{k,h}(t,y)=\sum_{w,x}
  K(t,y;0,w)S(w,x)_{hh}\eta_k(x)
  \enum
}
for a representative ensemble of gauge fields,
all $k=1,\ldots,N_s$ and $h=r,s$.

For a given gauge field,
the computation of $\phi_{k,h}(t,y)$ proceeds in two steps,
where one first calculates $\phi_{k,h}(0,w)$ by solving
the lattice Dirac equation with mass $m_{0,h}$ and
source field $\eta_k(x)$. The calculated field must then 
be evolved in flow time from time $0$ to time $t$.
In the direction of increasing flow time,
the flow equations
are numerically stable and the solution can easily 
be obtained, with negligible
integration errors, using a Runge--Kutta integrator
(appendix D). 


\lussubsec 7.2 Chiral condensate

In the case of the time-dependent quark condensate,
\equation{
  \langle S_t^{rr}(x)\rangle=
  -\sum_{v,w}\langle
  \tr\{K(t,x;0,v)\left[S(v,w)_{rr}-\cfl\delta_{vw}\right]
  K(t,x;0,w)^{\dagger}\}\rangle,
  \enum
}
an averaging over the position $x$ is again desirable. 
Introducing random source fields as above, one is then
led to the representation
\equation{
  {1\over N_{\Gamma}}\sum_{x\in\Gamma}
  \langle S_t^{rr}(x)\rangle=
  -{1\over N_{\Gamma}}
  \sum_{v,w}\langle
  \xi_k(t;0,v)^{\dagger}\left[S(v,w)_{rr}-\cfl\delta_{vw}\right]
  \xi_k(t;0,w)\rangle
  \enum
}
in terms of the fields
\equation{
  \xi_k(t;s,w)=\sum_x K(t;x;s,w)^{\dagger}\eta_k(x)
  \enum
}
at flow time $s=0$. 

The computation of these fields requires 
the {\it adjoint flow equation}
\equation{
  (\partial_s+\Delta)\xi_k(t;s,w)=0
  \enum
}
to be solved from time $s=t$ (where $\xi_k$ coincides
with $\eta_k$) to $s=0$. A straightforward Runge--Kutta
integration should not be used here, because
the Laplacian $\Delta$ is the one in presence
of the gauge field determined by the gradient flow at time $s$,
which would therefore have to be evolved backward in time,
i.e.~in the unstable direction.
As explained in appendix E, this difficulty can be overcome
through a hierarchical scheme that avoids the backward integration
of the gauge field.

